\newpage

\chapter{Fundamentação Teórica}

\section{Solo}

%2.1 SOLO - nesse subitem você pode definir o solo (produto do intemperismo sobre as rochas) e fazer comentários sobre sua origem (solos residuais, solos transportados e as características dos horizontes), fazer comentários sobre o tamanho das partículas, comentários sobre o comportamento desse material com a variação da umidade (limites de consistência)

O solo constitui um material natural da crosta terrestre gerado a partir da decomposição ou da ruptura mecânica, química e bioquímica das rochas subjacentes, por meio de um processo contínuo conhecido como intemperismo, operante na interface entre a litosfera e a atmosfera \cite{Kovacs2017, Barroso2019, Junior2021, Soler2021}. Quanto à sua origem e formação geológica, os solos classificam-se fundamentalmente em dois grupos: os solos residuais, formados \textit{in situ} e que permanecem sobrejacentes à rocha matriz original, e os solos transportados, constituídos por partículas que sofreram erosão e foram deslocadas de seu local de origem por agentes como água, vento ou gravidade \cite{Junior2021, Silva2019, Barroso2019, Soler2021}. O avanço desse perfil intempérico dita a variação das propriedades físicas ao longo da profundidade \cite{Rahardjoa2004}. Com o tempo e os diferentes graus de exposição ambiental, o terreno desenvolve camadas sobrepostas denominadas horizontes pedológicos, os quais apresentam características nítidas, exibindo horizontes superiores tipicamente mais finos e ricos em umidade, que evoluem verticalmente até o substrato rochoso alterado ou em processo de decomposição \cite{Rahardjoa2004, Tangari2021}.

Fisicamente, a distribuição granulométrica (tamanho das partículas) funciona como a base quantitativa primordial para distinguir e antecipar o comportamento mecânico e hidráulico dos solos \cite{Castro2023, Liu2024}. Essa caracterização divide o material em frações grossas (como pedregulhos e areias) e frações finas (siltes e argilas), sendo que, para os solos finos, a plasticidade e a mineralogia ditam as propriedades gerais de forma muito mais eficaz do que a análise textural isolada \cite{MorenoMaroto2021}. Desse modo, o comportamento mecânico da matriz fina sofre severa influência das variações em seu teor de umidade \cite{Demir2021}. Para quantificar essas transições, utilizam-se os limites de consistência de Atterberg, nos quais o \gls{lc} determina a transição do estado sólido para o semissólido, ponto abaixo do qual reduções na umidade não alteram o volume do material; o \gls{lp} separa os estados semissólido e plástico, representando o limiar em que o solo pode ser moldado sem se fraturar; e o \gls{ll} demarca a transição para o estado de comportamento fluido e de baixa resistência ao cisalhamento \cite{Demir2021, MorenoMaroto2021}. A diferença matemática entre o \gls{ll} e o \gls{lp} origina o \gls{ip}, parâmetro que reflete a amplitude hídrica em que o solo exibe plasticidade, correlacionando-se diretamente com o potencial de contração, expansão e sensibilidade à umidade \cite{Smith1985}.

\section{Solo como material construtivo}

%2.2 SOLO COMO MATERIAL CONSTRUTIVO - nesse subitem você pode retomar os benefícios dessas técnicas construtivas, tanto em aspectos ambientais (disponibilidade do material, reciclabilidade, consumo de energia, geração de resíduos) como em relação ao conforto dos usuários (térmico, acústico, etc.). Aqui você pode citar brevemente algumas técnicas como taipa de pilão, cob, pau a pique, bloco de terra comprimida e o próprio adobe. 

O uso da terra na construção civil tem voltado a ganhar destaque nas últimas décadas. Esse crescimento está relacionado à busca por soluções que reduzam os impactos ambientais, diminuam o consumo de recursos naturais e contribuam para práticas mais sustentáveis. Nesse contexto, as construções em terra se destacam por aproveitar materiais disponíveis localmente, gerar menos impactos ao meio ambiente e unir conhecimentos tradicionais com técnicas e tecnologias atuais.

%Benefícios Ambientais

A principal vantagem ambiental da terra está no fato de ser um material abundante e facilmente encontrado em diferentes regiões, o que reduz a necessidade de transporte e, consequentemente, os impactos gerados por essa etapa. Além disso, diferentemente de materiais como tijolos cerâmicos, a terra pode ser utilizada sem a necessidade de queima em altas temperaturas, o que diminui o consumo de energia e a emissão de gases poluentes durante sua produção.

Nesse sentido, \citeonline{Zhang2025} apontam que os materiais à base de terra consomem, em média, apenas 6\% da energia utilizada na produção do concreto e geram cerca de 3\% das emissões de carbono associadas aos tijolos cerâmicos convencionais. De forma semelhante, \citeonline{Fernandes2019}, por meio de uma análise de \gls{lca}, verificaram que o \gls{btc} apresenta uma energia incorporada de apenas 3,94 MJ por bloco e um \gls{gwp} de 0,39 kg $CO_2$ eq por bloco. Esses resultados indicam uma redução de aproximadamente 50\% nos impactos ambientais potenciais quando comparado aos sistemas convencionais de alvenaria.

Em relação ao fim da vida útil, a terra crua sem aditivos químicos se destaca por seu potencial de reutilização e reintegração ao meio ambiente. Dentro da lógica do conceito "do berço ao berço" (\textit{cradle-to-cradle}), o material pode ser destorroado, reumedecido e reutilizado diversas vezes. Segundo \citeonline{Morel2021, Aspillaga2024}, a terra é 100\% reciclável e biodegradável, podendo retornar ao ambiente natural sem gerar resíduos poluentes.
 
\begin{figure}[p]
    \centering
    \caption{\label{fig_diagrama_fases}Diagrama de entradas e saídas dos processos para construir uma parede de taipa de pilão.}
 
    \resizebox{0.97\textwidth}{!}{%
    \begin{tikzpicture}[
        >={Stealth[length=3.5mm]},
        % ---------- BLOCOS ----------
        bloco_laranja/.style={
            rectangle, fill=orange!45,
            minimum width=4.4cm, text width=4.1cm, minimum height=1.2cm,
            text centered, rounded corners=4pt,
            font=\sffamily\small\bfseries},
        bloco_azul/.style={
            rectangle, fill=cyan!75!black,
            minimum width=4.4cm, text width=4.1cm, minimum height=1.2cm,
            text centered, rounded corners=4pt,
            font=\sffamily\small\bfseries},
        bloco_roxo/.style={
            rectangle, fill=violet!55,
            minimum width=4.4cm, text width=4.1cm, minimum height=1.2cm,
            text centered, rounded corners=4pt,
            font=\sffamily\small\bfseries},
        bloco_cinza/.style={
            rectangle, fill=gray!55,
            minimum width=4.4cm, text width=4.1cm, minimum height=1.2cm,
            text centered, rounded corners=4pt,
            font=\sffamily\small\bfseries},
        bloco_verde/.style={
            rectangle, fill=green!55,
            minimum width=4.4cm, text width=4.1cm, minimum height=1.2cm,
            text centered, rounded corners=4pt,
            font=\sffamily\small\bfseries},
        % ---------- TEXTOS ----------
        txt_in/.style={font=\sffamily\footnotesize, align=right, text width=2.4cm},
        txt_out/.style={font=\sffamily\footnotesize, align=left,  text width=3.2cm},
        txt_label/.style={font=\sffamily\footnotesize, align=left, text width=2.7cm},
        txt_fase/.style={font=\sffamily\small\bfseries, rotate=90, align=center},
        txt_hdr/.style={font=\sffamily\small\bfseries},
        % ---------- SETAS ----------
        seta/.style={->,  draw=gray!65, line width=3.2pt},
        seta_loop/.style={->, draw=gray!55, line width=1.4pt},
    ]
 
    % ==============================================
    % LAYOUT DE COLUNAS (coordenadas X)
    %   col_fase_txt = -8.5   texto vertical fase
    %   col_lim_esq  = -7.5   limite esq. das linhas
    %   col_div_vert = -6.2   divisória vertical
    %   col_lbl      = -6.1   início rótulos (A1, A4…)
    %   col_in_east  = -0.3   âncora east das entradas
    %   col_proc     =  3.2   centro dos processos
    %   col_out_west =  6.6   âncora west das saídas
    %   col_lim_dir  = 10.5   limite dir. das linhas
    % ==============================================
 
    % ==============================================
    % TÍTULO
    % ==============================================
    \node[font=\sffamily\bfseries\normalsize] at (3.2, 2) {TAIPA DE PILÃO};
 
    % ==============================================
    % CABEÇALHO
    % ==============================================
    \node[txt_hdr] at ( -1.5, 1.0) {ENTRADAS};
    \node[txt_hdr] at ( 3.2, 1.0) {PROCESSOS};
    \node[txt_hdr] at ( 7.5, 1.0) {SAÍDAS};
 
    \draw[thick] (-6, 0.5) -- (10.5, 0.5);
 
    % ==============================================
    % DIVISÓRIA VERTICAL DE FASES
    % ==============================================
    %\draw[thick] (-6.2, 0.4) -- (-6.2, -21.0);
 
    % ==============================================
    % RÓTULOS VERTICAIS DAS FASES
    % ==============================================
    \node[txt_fase] at (-7.0, -0.5)  {Produto};
    \node[txt_fase] at (-7.0, -6.5)  {Construção};
    \node[txt_fase] at (-7.0, -14.5)  {Uso};
    \node[txt_fase] at (-7.0, -19.5) {Fim da Vida Útil};
 
    % ==============================================
    % A1 — RAW MATERIALS SUPPLY  y = -1.0
    % ==============================================
    \node[txt_label, anchor=north west] at (-6.1,  0)
        {\textbf{A1}\\[-2pt]Fornecimento de matéria-prima};
 
    \node[txt_in,  anchor=east] (inA1) at (-0.3, -0.5) {Combustível};
    \node[bloco_laranja]        (prA1) at ( 3.2, -0.5)
        {Fornecimento de matérias-primas\\(extração)};
    \node[txt_out, anchor=west] (ouA1) at ( 6.6, -0.5)
        {Solo orgânico e rochas\\Emissões para o ar};
 
    \draw[seta] (inA1.east) -- (prA1.west);
    \draw[seta] (prA1.east) -- (ouA1.west);
 
    % Separador Product / A4
    \draw[thick] (-6, -1.5) -- (10.5, -1.5);
 
    % ==============================================
    % A4 — TRANSPORT OF MATERIALS  y = -3.1
    % ==============================================
    \node[txt_label, anchor=north west] at (-6.1, -2)
        {\textbf{A4}\\[-2pt]Transporte};
 
    \node[txt_in,  anchor=east] (inA4) at (-0.3, -2.4) {Combustível};
    \node[bloco_azul]           (prA4) at ( 3.2, -2.4)
        {Transporte de materiais e equipamentos};
    \node[txt_out, anchor=west] (ouA4) at ( 6.6, -2.4) {Emissões para o ar};
 
    \draw[seta] (inA4.east) -- (prA4.west);
    \draw[seta] (prA4.east) -- (ouA4.west);
    \draw[seta] (prA1.south) -- (prA4.north);
 
    % Separador leve A4 / Construção
    \draw[gray!40, thin] (-6, -3.25) -- (10.5, -3.3);
 
    % ==============================================
    % A5 — CONSTRUCTION
    % ==============================================
    \node[txt_label, anchor=north west] at (-6.1, -7)
        {\textbf{A5}\\[-2pt]Instalação/\\Construção};
 
    % --- Formwork assembly  y = -5.1 ---
    \node[txt_in, anchor=east] (inFA) at (-0.3, -4.2) {Solo\\Cal};
    \node[bloco_azul]          (prFA) at ( 3.2, -4.2) {Montagem de fôrmas};
 
    \draw[seta] (inFA.east) -- (prFA.west);
    \draw[seta] (prA4.south) -- (prFA.north);
 
    % --- Mixing  y = -6.6 ---
    \node[txt_in,  anchor=east] (inMX) at (-0.3, -6)
        {Solo\\Cal\\Água\\Combustível};
    \node[bloco_azul]           (prMX) at ( 3.2, -6) {Mistura};
    \node[txt_out, anchor=west] (ouMX) at ( 6.6, -6)
        {Resíduos de embalagens de cal\\Emissões para o ar};
 
    \draw[seta] (inMX.east) -- (prMX.west);
    \draw[seta] (prMX.east) -- (ouMX.west);

    \draw[seta] (prFA.south) -- (prMX.north);
    % --- Filling the formwork  y = -8.1 ---
    \node[txt_in,  anchor=east] (inFF) at (-0.3, -7.8) {Combustível};
    \node[bloco_azul]           (prFF) at ( 3.2, -7.8) {Preenchimento da fôrma};
    \node[txt_out, anchor=west] (ouFF) at ( 6.6, -7.8)
        {Resíduos de material misto\\Emissões para o ar};
 
    \draw[seta] (inFF.east) -- (prFF.west);
    \draw[seta] (prFF.east) -- (ouFF.west);
    \draw[seta] (prMX.south) -- (prFF.north);
 
    % Loop retorno (base de Mixing de volta para as entradas)
    \draw[seta_loop]
        ($(ouFF)+(0,0.7)$) -- ($(ouFF.north west)+(1.72,0.1)$) -- ($(prFF.north west)+(-1,0.2)$) -- ($(prMX.south west)+(-1,0.2)$) -- ($(prMX.south west)+(0,0.2)$);
 
    % --- Ramming  y = -9.5 ---
    \node[txt_in,  anchor=east] (inRM) at (-0.3, -9.7) {Combustível};
    \node[bloco_azul]           (prRM) at ( 3.2, -9.7) {Apiloamento};
    \node[txt_out, anchor=west] (ouRM) at ( 6.6, -9.7) {Emissões para o ar};
 
    \draw[seta] (inRM.east) -- (prRM.west);
    \draw[seta] (prRM.east) -- (ouRM.west);
    \draw[seta] (prFF.south) -- (prRM.north);
 
    % --- Formwork disassembly  y = -10.8 ---
    \node[bloco_azul] (prFD) at (3.2, -11.6) {Desmontagem da fôrma};
    \draw[seta] (prRM.south) -- (prFD.north);
 
    % Caixa tracejada verde — A5
    %\draw[dashed, green!55!black, line width=1.2pt, rounded %corners=5pt]
    %    (-4.5, -4.45) rectangle (5.7, -11.45);
 
    % Separador Construção / Use
    \draw[thick] (-6, -12.5) -- (10.5, -12.5);
 
    % ==============================================
    % B2 — USE / MAINTENANCE
    % ==============================================
    \node[txt_label, anchor=north west] at (-6.1, -14)
        {\textbf{B2}\\[-2pt]Manutenção};
 
    % --- Transport (Use)  y = -13.1 ---
    \node[txt_in,  anchor=east] (inTU) at (-0.3, -13.5) {Combustível};
    \node[bloco_roxo]           (prTU) at ( 3.2, -13.5) {Transporte};
    \node[txt_out, anchor=west] (ouTU) at ( 6.6, -13.5) {Emissões para o ar};
 
    \draw[seta] (inTU.east) -- (prTU.west);
    \draw[seta] (prTU.east) -- (ouTU.west);
    \draw[seta] (prFD.south) -- (prTU.north);
 
    % --- Whitewashing  y = -14.5 ---
    \node[txt_in,  anchor=east] (inWS) at (-0.3, -15.5) {Água de\\cal};
    \node[bloco_roxo]           (prWS) at ( 3.2, -15.5)
        {Caiação/\\tratamento de superfície};
    \node[txt_out, anchor=west] (ouWS) at ( 6.6, -15.5)
        {Resíduos de embalagens de cal};
 
    \draw[seta] (inWS.east) -- (prWS.west);
    \draw[seta] (prWS.east) -- (ouWS.west);
    \draw[seta] (prTU.south) -- (prWS.north);
 
    % Separador Use / End-of-life
    \draw[thick] (-6, -16.7) -- (10.5, -16.7);
 
    % ==============================================
    % END-OF-LIFE
    % ==============================================
 
    % --- C1 Demolition  y = -16.55 ---
    \node[txt_label, anchor=north west] at (-6.1, -17.2)
        {\textbf{C1}\\[-2pt]Demolição};
 
    \node[txt_in,  anchor=east] (inDM) at (-0.3, -17.7) {Combustível};
    \node[bloco_cinza]          (prDM) at ( 3.2, -17.7) {Demolição};
    \node[txt_out, anchor=west] (ouDM) at ( 6.6, -17.7) {Emissões para o ar};
 
    \draw[seta] (inDM.east) -- (prDM.west);
    \draw[seta] (prDM.east) -- (ouDM.west);
    \draw[seta] (prWS.south) -- (prDM.north);
 
    % Separador C1 / C3
    \draw[gray!40, thin] (-6, -18.5) -- (10.5, -18.5);
 
    % --- C3 Milling/Sieving  y = -18.1 ---
    \node[txt_label, anchor=north west] at (-6.1, -18.7)
        {\textbf{C3}\\[-2pt]Processamento de resíduos};
 
    \node[txt_in,  anchor=east] (inML) at (-0.3, -19.5) {Eletricidade};
    \node[bloco_cinza]          (prML) at ( 3.2, -19.5) {Moagem/Peneiramento};
 
    \draw[seta] (inML.east) -- (prML.west);
    \draw[seta] (prDM.south) -- (prML.north);
 
    % Separador C3 / C4
    \draw[gray!40, thin] (-6, -20.3) -- (10.5, -20.3);
 
    % --- C4 Disposal on site  y = -19.5 ---
    \node[txt_label, anchor=north west] at (-6.1, -20.7)
        {\textbf{C4}\\[-2pt]Descarte};
 
    \node[bloco_cinza] (prDS) at (5.4, -21.2) {Descarte no local};

    \draw[seta, dashed] (prDM.south east) -- (prDS);
 
    % Separador C4 / D
    \draw[thick] (-6, -22.1) -- (10.5, -22.1);
 
    % --- D Recovery/Recycling  y = -20.85 ---
    \node[txt_label, anchor=north west] at (-6.1, -22.2)
        {\textbf{D}\\[-2pt]Reutilização, Recuperação\\Reciclagem};
 
    \node[bloco_verde] (prRR) at (3.2, -23) {Recuperação / Reciclagem};
    \draw[seta] (prML.south west) -- (prRR.north west);
 
    % Seta triângulo verde Formwork → Mixing
    \draw[->, draw=green!55!black, line width=3.2pt,
          >={Triangle[length=4mm,width=2mm]}]
        (prFA.south west |- -5,-5.3) -- (-5,-5.5 -| prMX.north west);
 
    % Linha tracejada verde larga (contorno A4 → D, lado esquerdo)
    \draw[dashed, green!55!black, line width=1.2pt]
        (prMX.north west)
        -- (prFA.south west |- prFA.south west)
        -- (-2.5, -4.8 |- prFA.south west)
        -- (-2.5, -23)
        -- (prRR.west |- 0,-23)
        -- (prRR.west);
 
    % Linha final
    \draw[thick] (-6, -24) -- (10.5, -24);
 
    \end{tikzpicture}%
    }% fim do resizebox
 
    \vspace{0.2cm}
    \fonte{\citeonline[p. 9, adaptado pelo autor]{Fernandes2019}.}
\end{figure}


%Ressalvas Metodológicas: O Impasse da Estabilização Química

Apesar das vantagens ambientais da terra como material de construção, a literatura aponta um desafio relacionado ao uso de estabilizantes químicos. Para reduzir a degradação da terra crua quando exposta à água, a adição de cimento Portland tornou-se uma prática comum. No entanto, \citeonline{MoraRuiz2025, Fernandes2019} destacam que a incorporação desse material aumenta significativamente o \gls{gwp} dos compósitos. Como consequência, parte dos benefícios ambientais da terra pode ser perdida, reduzindo sua vantagem em relação a outros materiais construtivos. Esse cenário reforça a importância de pesquisas voltadas ao desenvolvimento de estabilizantes alternativos, especialmente aqueles produzidos a partir de resíduos ou materiais de origem orgânica.

%Conforto dos Usuários

Além dos benefícios ambientais, a terra também se destaca por contribuir para o conforto térmico e acústico das edificações, favorecendo melhores condições de habitabilidade.

No aspecto térmico, os sistemas construtivos em terra apresentam alta capacidade de absorver e liberar umidade, além de elevada inércia térmica. Segundo \citeonline{Giada2019}, o solo funciona de maneira semelhante a um \gls{pcm}. Seus poros permitem a evaporação da água, absorvendo calor do ambiente em períodos quentes, e sua posterior condensação em períodos mais frios, liberando calor. Esse processo ajuda a regular passivamente a temperatura interna e a manter a umidade relativa do ar em níveis considerados saudáveis, entre 40\% e 60\%.

Os efeitos dessas propriedades podem ser observados em estudos como o de \citeonline{Rincon2019}, realizado em Burkina Faso. Os autores verificaram que o uso da terra, combinado com estratégias de sombreamento e ventilação noturna, reduziu o desconforto térmico anual para apenas 209 horas. Além disso, as edificações apresentaram elevada estabilidade térmica, com variações internas de apenas 2ºC, mesmo diante de oscilações externas muito mais intensas.

Em relação ao desempenho acústico, \citeonline{PachecoTorgal2012} destacam que a elevada massa e espessura das paredes de terra crua proporcionam bom isolamento sonoro. Esses sistemas podem atingir  \gls{sri} entre 46 dB e 57 dB, valores superiores aos normalmente observados em alvenarias de tijolos cerâmicos convencionais, que apresentam desempenho em torno de 35 dB.

%Limitações Físicas e a Necessidade de Desenho Bioclimático

Apesar dos resultados positivos, o desempenho térmico da terra não deve ser valiado de maneira simplista. \citeonline{Kyriakidis2018} destacam que, embora as paredes de adobe proporcionem conforto térmico devido à sua elevada massa térmica, elas também apresentam alto valor-U (transmitância térmica). Isso indica que a terra permite a transferência de calor com maior facilidade quando comparada aos materiais isolantes modernos.

Da mesma forma, \citeonline{Rincon2019} ressaltam que, em regiões de clima extremo, os benefícios da alta inércia térmica podem ser reduzidos na ausência de estratégias bioclimáticas adequadas. Sem medidas como o isolamento da cobertura para minimizar a incidência de radiação solar, as paredes de terra podem absorver calor em excesso durante o dia e liberá-lo para o ambiente interno em horários indesejados, causando desconforto térmico aos ocupantes.

%Técnicas Construtivas de Terra

Segundo \citeonline{Arduin2022, Silva2019, Morel2021, MoraRuiz2025}, a terra pode ser utilizada na construção por meio de diversas técnicas, desde métodos tradicionais até tecnologias mais recentes. As principais diferenças entre essas técnicas estão relacionadas à umidade do solo, ao processo de preparação do material e à forma como ele é aplicado na estrutura:

\begin{alineas}
    \item adobe: é a técnica mais antiga e simples de construção com terra. Ela consiste na produção de blocos a partir de uma mistura de solo e água em estado plástico, geralmente acrescida de fibras vegetais. Após a moldagem, os blocos são secos naturalmente ao sol, sem a necessidade de processos de compactação mecânica \cite{Arduin2022, MoraRuiz2025}.
    \item bloco de terra comprimida (BTC): é considerado uma evolução tecnológica do adobe, o \gls{btc} é produzido com solo levemente úmido submetido à compactação mecânica em fôrmas, utilizando prensas manuais ou hidráulicas. Esse processo resulta em blocos mais densos, com dimensões mais uniformes e maior resistência mecânica \cite{Arduin2022}.
    \item taipa de pilão: consiste na construção de paredes monolíticas com função estrutural. A técnica baseia-se na compactação sucessiva de camadas de solo úmido dentro de fôrmas verticais removíveis, que podem ser de madeira ou metal. O apiloamento pode ser realizado de forma manual ou mecanizada \cite{Silva2019, Arduin2022}.
    \item cob: é uma técnica construtiva monolítica executada diretamente no local da obra, sem a utilização de fôrmas. O material é composto por uma mistura de solo, água e palha, que é aplicada e moldada manualmente de forma gradual até a formação das paredes \cite{Arduin2022, Morel2021}.
    \item pau a pique (taipa de mão): é um sistema de vedação que utiliza uma estrutura interna de madeira ou bambu como suporte. Sobre essa trama, aplica-se manualmente uma mistura de barro em estado plástico, responsável pelo preenchimento e acabamento das paredes \cite{Silva2019, MoraRuiz2025}.
\end{alineas}

%Comparações Tecnológicas e Operacionais das Técnicas

Ao comparar as diferentes técnicas construtivas em terra, \citeonline{MoraRuiz2025, Arduin2022} apontam diferenças importantes em relação ao desempenho e à execução. Técnicas baseadas na moldagem manual, como o adobe e o cob, tendem a apresentar maior variabilidade na resistência mecânica quando utilizadas em seu estado natural. Além disso, são mais suscetíveis à retração e aos processos de degradação causados pela ação da água.

Por outro lado, técnicas que utilizam compactação, como o \gls{btc} e a taipa de pilão, resultam em elementos com maior densidade e melhor estabilidade estrutural. No entanto, essas vantagens dependem do uso de equipamentos e fôrmas específicos, como prensas manuais ou hidráulicas e sistemas de confinamento. Segundo \citeonline{MoraRuiz2025}, essa necessidade pode aumentar os custos de implantação, especialmente em regiões com infraestrutura técnica limitada.

\subsection{Adobe}

%2.2.1. Adobe - nesse subitem você pode se aprofundar um pouco mais na técnica do adobe: origem, regiões do planeta em que sua utilização é mais comum, tipo de solo que melhor se adapta à técnica, adições comuns para o melhoramento de suas propriedades (pelos e gordura animal, adição de fibras vegetais, etc) e estratégias para a durabilidade das construções (beirais, impermeabilização das fundações), etc. Nesse subitem foque em práticas tradicionais. O uso de produtos industrializados como cal e cimento serão abordados no próximo subitem.

\section{Técnicas de melhoramento de solos}

%2.3 TÉCNICAS DE MELHORAMENTO DE SOLOS - nesse subitem você pode contextualizar a necessidade de melhoramento de solos para que ele usado para fins construtivos, você pode trazer os principais tipos de estabilização: mecânica, granulométrica/física e química.

\subsection{Solo-cimento}

%2.3.1 Solo-cimento - nesse subitem você pode trazer mais detalhes sobre o solo-cimento, principais características, tipos de cimento comumente utilizados para esse fim, principais reações químicas, tipos de solo que melhor se adaptam à técnica e principais limitações (tendência de trincamento, perda de resistência quando exposto à umidade excessiva, etc.).












% A Fundamentação Teórica ou Revisão da Literatura deve contextualizar teoricamente o problema e apresentar o estágio atual de conhecimento sobre o assunto. A revisão de literatura é a parte que sustenta todo o projeto, pois é por meio dela que se encontram indicações para a solução do problema identificado. Devido a sua importância, ela deve ser atual, abrangente e com profundidade. Por outro lado, deve deter-se a assuntos específicos do projeto de pesquisa em desenvolvimento, sem incluir aspectos desnecessários e de forma exaustiva. Trata-se da apresentação do embasamento teórico sobre o qual se fundamentará o trabalho. Durante a escrita, o autor deve lembrar de que o texto não deve ser construído apenas com uma sequência de recortes e citações de outros autores sobre o tema. Segundo Gil (2018), a revisão bibliográfica não é constituída apenas por referências ou sínteses do relato de estudos, mas por discussão crítica das obras citadas. 

% ORIENTAÇÕES GERAIS Escreva todo seu texto utilizando estilo normal, fonte Times New Roman, tamanho 12, exceto as citações com mais de três linhas, notas de rodapé, paginação, legendas e fontes das figuras e tabelas, que devem ser em tamanho menor e uniforme. Todo os elementos prétextuais e textuais deverão ser formatados para um tamanho de página A4 (210 x 297 mm), limitado por margens esquerda e superior de 3 cm e direita e inferior de 2 cm. Todo o texto deve ser digitado com espaçamento 1,5 entre as linhas e alinhamento justificado, exceto as citações de mais de três linhas, notas de rodapé, referências, legendas das figuras e tabelas. As referências bibliográficas, ao final do trabalho, devem ser separadas entre si por um espaço simples em branco. Os títulos das seções primárias devem começar em página ímpar (anverso) e ser separados do texto que os sucede por um espaço entre as linhas de 1,5. Do mesmo modo, os títulos das subseções devem ser separados do texto que os procede e que os sucede por um espaço entre as linhas de 1,5. Lista de figuras, lista de tabelas, lista de abreviaturas e siglas, lista de símbolos, sumário, referências, apêndice(s) e anexo(s) são títulos sem indicativo numérico e devem ser centralizados. 287 As folhas ou páginas pré-textuais devem ser contadas, mas não numeradas. Os trabalhos devem ser digitados ou datilografados somente no anverso. Todas as folhas, a partir da folha de rosto, devem ser contadas sequencialmente, considerando somente o anverso (a capa não é contada - página “zero”). A numeração deve aparecer, a partir da primeira folha da parte textual (Introdução), em algarismos arábicos, no canto superior direito da folha. 

% Tabelas As tabelas devem ser enumeradas sequencialmente, citadas no texto, inseridas o mais próximo possível do trecho a que se referem e padronizadas conforme o Instituto Brasileiro de Geografia e Estatística (IBGE, 1993). • Toda tabela que ultrapassar a dimensão da página em número de linhas e tiver poucas colunas, pode ter o centro apresentado em duas ou mais partes, lado a lado, na mesma página, separando-se as partes por um traço vertical duplo e repetindo-se o cabeçalho; • Toda tabela que ultrapassar a dimensão da página em número de colunas, e tiver poucas linhas, pode ter o centro apresentado em duas ou mais partes, uma abaixo da outra, na mesma página, repetindo-se o cabeçalho das colunas indicadoras e os indicadores de linha; • Para toda tabela que ultrapassar as dimensões da página: cada página deve ter o conteúdo do topo e o cabeçalho da tabela ou o cabeçalho da parte; cada página deve ter uma das seguintes indicações: continua para a primeira, conclusão para a última e continuação para as demais. O título das tabelas deve ser incluído na linha imediatamente anterior à tabela e centralizado, utilizando fonte Times New Roman, tamanho 10 e cor preta, conforme Tabela 1. Para configurar o título das tabelas acesse o menu “Referências” e clique em “Inserir Legenda”. A fonte deve ser identificada logo abaixo, mesmo que de autoria própria, indicando “acervo particular” ou “autoria própria”.


% Equações e fórmulas Buscando facilitar a leitura, as equações e fórmulas devem ser destacadas no texto e, se necessário, numeradas com algarismos arábicos entre parênteses, alinhados à direita. Todas as variáveis envolvidas nas equações bem como a unidade do parâmetro calculado devem ser explicitadas ao longo do texto ou em seguida à apresentação da equação. A Equação 1 apresenta a fórmula para o cálculo da velocidade escalar média, como forma de exemplificação. vm = ∆S ∆T (1) Onde: vm: velocidade escalar média (m/s); ΔS: distância total (m); ΔT: intervalo de tempo (s). 

% Ilustrações Qualquer que seja o tipo de ilustração, sua identificação aparece na parte superior, precedida da palavra designativa (desenho, fluxograma, gráfico, mapa, planta, figura, imagem, entre outros), seguida de seu número de ordem de ocorrência no texto, em algarismos arábicos, e do respectivo título. A fonte deve ser identificada logo abaixo, mesmo que de autoria própria, indicando “acervo particular” ou “autoria própria”. A ilustração deve ser citada no texto e inserida o mais próximo possível do trecho a que se refere. A Figura 1 mostra uma representação esquemática de um solo com estrutura alveolar, com a intenção de demonstrar o processo de configuração de figuras.

%\begin{figure}[H]
%    \centering
%    \caption{Teste}
%    \label{fig:placeholder}
%    \includegraphics[width=0.5\linewidth]{image.jpg}
%    
%    {\small Fonte: Elaborado pelo autor (2026). \par}
%\end{figure}

% Citações As referências a autores ou transcrição de informações retiradas de outras fontes devem seguir as diretrizes da ABNT NBR 10520. As citações diretas consistem na transcrição textual de parte da obra do autor consultado, e as citações indiretas consistem na elaboração de um texto baseado na obra do autor consultado. Nas citações, as chamadas pelo sobrenome do autor devem ser em letras maiúsculas e minúsculas e, quando estiverem entre parênteses, devem ser em letras maiúsculas. Exemplo 1: De acordo com Miguez, Veról e Rezende (2015), o processo de urbanização gera grandes modificações no ambiente natural, alterando os padrões de uso e ocupação do solo e agravando os problemas de enchentes. Exemplo 2: A mistura entre agregado reciclado cimentício e agregado reciclado de cerâmica vermelha pode favorecer a formação de compostos cimentantes, uma vez que a fração fina do resíduo de cerâmica vermelha tem desempenho de material pozolânico (SILVA, 2014). As citações diretas, no texto, de até três linhas, devem estar contidas entre aspas duplas. As citações diretas, no texto, com mais de três linhas, devem ser destacadas com recuo de 4 cm da margem esquerda, com letra menor que a do texto utilizado e sem as aspas. Nas citações diretas deve-se especificar no texto: a(s) página(s), volume(s), tomo(s) ou seção(ões) da fonte consultada. Estas informações devem seguir o ano da publicação, separado(s) por vírgula e precedido(s) pelo termo que o(s) caracteriza, de forma abreviada. Nas citações indiretas, a indicação da(s) página(s) consultada(s) não é obrigatória. Exemplo de citação direta de até três linhas: “Cimento, no sentido geral da palavra pode ser descrito como um material com propriedades adesivas e coesivas que o fazem capaz de unir fragmentos minerais na forma de uma unidade compacta.” (NEVILLE, 2015, p.1). Exemplo de citação direta com mais de três linhas: Uma geogrelha é definida como um material polimérico (isto é, geossintético) composto por conjuntos paralelos de arestas elásticas conectadas com aberturas do tamanho suficiente para permitir a penetração do solo, pedras ou outros materiais geotécnicos do entorno. As geogrelhas normalmente são feitas com polietileno de alta densidade (HDPE) e polipropileno (PP). A principal função de uma geogrelha é o reforço. (DAS; SOBHAN, 2019, p. 662).